Las arquitecturas informadas por física han demostrado ser una alternativa relevante para aproximar soluciones de ecuaciones diferenciales parciales (PDEs, por sus siglas en inglés) en escenarios donde los métodos clásicos pueden presentar limitaciones asociadas al costo computacional, la formulación de problemas inversos o la inferencia rápida. Sin embargo, los métodos actuales todavía presentan dificultades para identificar y priorizar regiones del dominio donde la representación aprendida resulta insuficiente, especialmente en problemas con comportamiento oscilatorio, múltiples escalas o acoplamiento físico fuerte.

Este trabajo explora el diseño e integración experimental de mecanismos de focalización adaptativa sobre arquitecturas de Redes de Kolmogorov-Arnold Informadas por Física (PIKANs, por sus siglas en inglés). Estos mecanismos buscan orientar el aprendizaje hacia regiones de mayor complejidad física durante el entrenamiento, utilizando señales como el residual físico o la variación local de la solución. La propuesta no demuestra la superioridad general de una arquitectura sobre otra, sino que caracteriza el comportamiento del mecanismo propuesto bajo condiciones controladas.

La validación sigue una estructura progresiva, desde problemas canónicos con solución analítica conocida o referencia numérica reproducible hasta sistemas físicos acoplados de mayor complejidad. El análisis considera métricas como el error relativo, el residual físico, la estabilidad del entrenamiento y el costo computacional por configuración. A partir de los resultados, el trabajo identifica condiciones de operación, limitaciones observadas y oportunidades de diseño para variantes PIKAN más robustas en contextos de física informada.

\vspace{0.8em}
\noindent\textbf{Palabras clave:} redes informadas por física, Kolmogorov-Arnold, focalización adaptativa, ecuaciones diferenciales parciales, aprendizaje automático científico, piezoelectricidad.